% Einheit 3 von 4 – Parameter aus Maßbedingungen (bank/scharen-von-geraden-und-ebenen/e3.jsonl, zusammenbau v0.1)
%% TODO Zeitmarke Einheit 3: Marken-Zeile nicht lesbar
%% TODO Zweigzeile Teil 1: was hier gelernt wird (Satz in Schülersprache aus dem Einheitstitel) – Einheit 3
\einheitenkopf[e3]{Einheit 3 von 4 · Parameter aus Maßbedingungen}
\zweigzeile{Abitur LK}
\setcounter{aufgabe}{13}

%% TODO Titel: Ich-kann-Satz fehlt in der Bank (Platzhalter: Kettenname) – Nr. 14
%% TODO Anweisung: ein Satz über der ersten Teilaufgabe fehlt in der Bank – Nr. 14
\begin{aufgabe}{Parameter aus Maßbedingungen}
%% TODO \rechenplatz{n}: Schrittzahl des Grundfalls fehlt in der Bank (nur bei fester Schreibform, 2.3 b)
\begin{teile}
\teil Für alle oder für eins? „Zeige, dass der Punkt $P$ von allen Ebenen $E_k$ denselben Abstand hat.“ \\ \kreuz{für alle Mitglieder} \\ \kreuz{Parameterwert gesucht}
\teil Für welche $k$ schließt $E_k: k x_1 + 3x_2 + 5x_3 = 2$ mit der $x_1x_2$-Ebene einen Winkel von $45^\circ$ ein? $k_1 = $ \leerfeld , $k_2 = $ \leerfeld
\teil Für welche $k$ schneidet die $x_3$-Achse die Ebene $E_k: k x_1 + x_2 + x_3 = 4$ unter $30^\circ$? $k_1 = $ \leerfeld , $k_2 = $ \leerfeld
\teil Für welche $k$ hat $P(2 | 1 | 0)$ von $E_k: x_1 + 2x_2 + 2x_3 = k$ den Abstand $2$? $k_1 = $ \leerfeld , $k_2 = $ \leerfeld
\teil Die Pyramide $ABCDS$ hat die Grundfläche $A(0 | 0 | 0)$, $B(8 | 0 | 0)$, $C(8 | 8 | 0)$, $D(0 | 8 | 0)$ und die Spitze $S(4 | 4 | 6)$. Jede Ebene $F_k: k x_1 + k x_2 + (k - 4) x_3 = 8k$ mit $0 < k < 4$ schneidet die Pyramide im Dreieck $BDQ_k$ mit $Q_k$ auf der Kante $SC$. Für ein $k$ ist der Flächeninhalt des Dreiecks am kleinsten. Bestimme dieses $k$. \hfill (Abitur 2025 LK) \\ $k = $ \leerfeld
\end{teile}
\end{aufgabe}

%% TODO Titel: Ich-kann-Satz fehlt in der Bank (Platzhalter: Kettenname) – Nr. 15
%% TODO Anweisung: ein Satz über der ersten Teilaufgabe fehlt in der Bank – Nr. 15
\begin{aufgabe}{Parameter aus Maßbedingungen – Fehler finden, Begründen}
\begin{teile}
\teil Finde den Fehler. Für welche $k$ schneidet die $x_3$-Achse die Ebene $E_k: k x_1 + x_2 + 3x_3 = 2$ unter $60^\circ$? \rechnung{\mathrm{cos}\,60^\circ &= \tfrac{3}{\sqrt{k^2 + 10}} \\ \sqrt{k^2 + 10} &= 6 \\ k^2 &= 26}
\teil Warum liefert die Bedingung „$P$ hat von $E_k$ einen vorgegebenen Abstand“ meist zwei Werte von $k$?
\end{teile}
\end{aufgabe}
